% requer booktabs. Diagnostico descritivo.
\begin{table}[htbp]\centering
\caption{Childcare enrolment and registered demand in the municipality of S\~ao Paulo, December snapshot, 2006--2017}
\label{tab:D15_E_anual_dez}
\small
\begin{tabular}{lrrrr}
\toprule
Year & Enrolment & Growth \% & Demand (queue) & Demand total \\
\midrule
2006 & 64,436 & -- & 129,594 & 129,594 \\
2007 & 83,364 & 29.4 & 79,231 & 79,231 \\
2008 & 109,717 & 31.6 & 57,607 & 61,089 \\
2009 & 123,155 & 12.2 & 74,707 & 74,707 \\
2010 & 130,412 & 5.9 & 100,401 & 100,401 \\
2011 & 195,561 & 50.0 & 101,165 & 101,165 \\
2012 & 214,094 & 9.5 & 97,160 & 97,160 \\
2013 & 214,460 & 0.2 & 96,696 & 96,696 \\
2014 & 228,204 & 6.4 & 94,191 & 94,191 \\
2015 & 260,785 & 14.3 & 75,194 & 75,194 \\
2016 & 284,179 & 9.0 & 65,040 & 65,040 \\
2017 & 296,260 & 4.3 & 44,092 & 44,092 \\
\bottomrule
\end{tabular}
\\[3pt]{\footnotesize \textit{Note:} Open data of the municipality (``Demanda Registrada e Matriculas'', historical series 2006--2017), summed over 96 districts, December of each year. MAT = childcare (creche) enrolment; DEM = registered childcare demand (the queue stock); DEM total adds preferential demand. The 2013 enrolment (214,460) matches the offered-places series used in the paper (VAGOFER 2013S1, 214,561). Source: build04\_diag15E.py; Portal de Dados Abertos PMSP.}
\end{table}
